%% ============================================================
%% project-book.tex
%% Compile with XeLaTeX in Overleaf:
%%   Menu -> Compiler -> XeLaTeX
%% ============================================================

\documentclass[12pt,a4paper]{article}

\usepackage{polyglossia}
\setmainlanguage{hebrew}
\setotherlanguage{english}

\usepackage{fontspec}
\newfontfamily\hebrewfont{FreeSans}
\newfontfamily\hebrewfonttt{FreeSans}
\setmainfont{FreeSans}

\usepackage{geometry}
\geometry{margin=2.5cm}

\usepackage{xcolor}
\usepackage{listings}
\usepackage{tikz}
\usetikzlibrary{shapes.geometric, arrows.meta, positioning, fit, backgrounds}
\usepackage{booktabs}
\usepackage{hyperref}
\usepackage{fancyhdr}
\usepackage{titlesec}
\usepackage{mdframed}
\usepackage{enumitem}
\usepackage{graphicx}

%% ── Colors ───────────────────────────────────────────────────
\definecolor{codebackground}{RGB}{245,245,245}
\definecolor{codekeyword}{RGB}{0,0,200}
\definecolor{codestring}{RGB}{163,21,21}
\definecolor{codecomment}{RGB}{0,128,0}
\definecolor{ubuntuorange}{RGB}{233,84,32}
\definecolor{fedorablue}{RGB}{60,110,180}
\definecolor{sectionblue}{RGB}{30,100,170}

%% ── Code listings ────────────────────────────────────────────
\lstset{
  backgroundcolor=\color{codebackground},
  basicstyle=\ttfamily\small,
  keywordstyle=\color{codekeyword}\bfseries,
  stringstyle=\color{codestring},
  commentstyle=\color{codecomment}\itshape,
  breaklines=true,
  frame=single,
  rulecolor=\color{gray!40},
  numbers=left,
  numberstyle=\tiny\color{gray},
  numbersep=8pt,
  showstringspaces=false,
  tabsize=4,
}

\lstdefinelanguage{bash}{
  keywords={if,then,else,fi,for,do,done,while,case,in,esac,function,return,
             set,source,export,local,echo,read,mapfile},
  morecomment=[l]{\#},
  morestring=[b]",
  morestring=[b]',
}

%% ── Page style ───────────────────────────────────────────────
\pagestyle{fancy}
\fancyhf{}
\fancyhead[R]{\small\rightmark}
\fancyhead[L]{\small multi-distro-provisioner}
\fancyfoot[C]{\thepage}
\renewcommand{\headrulewidth}{0.4pt}

%% ── Section formatting ───────────────────────────────────────
\titleformat{\section}
  {\Large\bfseries\color{sectionblue}}
  {\thesection}{1em}{}[\titlerule]

\titleformat{\subsection}
  {\large\bfseries}
  {\thesubsection}{1em}{}

%% ── Info box ─────────────────────────────────────────────────
\newmdenv[
  backgroundcolor=blue!5,
  linecolor=blue!40,
  linewidth=1.5pt,
  leftmargin=0pt,
  rightmargin=0pt,
  innerleftmargin=10pt,
  innerrightmargin=10pt,
  innertopmargin=8pt,
  innerbottommargin=8pt,
]{infobox}

%% ── TikZ styles ──────────────────────────────────────────────
\tikzstyle{block} = [rectangle, rounded corners, minimum width=3.5cm,
  minimum height=1cm, text centered, draw=black, fill=blue!10,
  font=\small]
\tikzstyle{decision} = [diamond, minimum width=3cm, minimum height=0.8cm,
  text centered, draw=black, fill=orange!20, font=\small, aspect=2]
\tikzstyle{arrow} = [thick,->,>=stealth]
\tikzstyle{startend} = [rectangle, rounded corners=8pt, minimum width=3cm,
  minimum height=0.9cm, text centered, draw=black, fill=green!20, font=\small]

%% ============================================================
\begin{document}
%% ============================================================

%% ── Title Page ───────────────────────────────────────────────
\begin{titlepage}
\centering
\vspace*{2cm}

{\Huge\bfseries\color{sectionblue} Multi-Distro Provisioner}\\[0.5cm]
{\Large מדריך מקיף לפרויקט אוטומציית תשתית לינוקס}\\[2cm]

{\large
\begin{tabular}{rl}
\textbf{מחברת:} & אסתי כהן \\[0.3cm]
\textbf{טכנולוגיות:} & Bash, Vagrant, systemd, nginx, Docker \\[0.3cm]
\textbf{הפצות:} & Ubuntu 22.04 \& Fedora 39 \\[0.3cm]
\end{tabular}
}

\vspace{2cm}

\begin{infobox}
\centering
\textit{"הדרך הטובה ביותר ללמוד לינוקס היא לבנות משהו שצריך לרוץ עליו."}
\end{infobox}

\vfill
{\small \today}
\end{titlepage}

\tableofcontents
\newpage

%% ============================================================
\section{הקדמה — למה בניתי את הפרויקט הזה}
%% ============================================================

במהלך לימודיי עבדתי אך ורק עם Ubuntu. זו הייתה ההפצה שהכרתי, זו שנלמדה בקורסים, זו שבה הרגשתי בנוח. אבל ככל שלמדתי יותר, הבנתי שבעולם האמיתי — במיוחד בסביבות ארגוניות — Ubuntu היא רק אחת מתוך אפשרויות רבות. RedHat, Fedora, CentOS ו-RHEL שולטות בשרתים ובענן הארגוני.

לכן שאלתי את עצמי: מה הדרך הטובה ביותר ללמוד בפועל כיצד הפצות לינוקס שונות עובדות? קריאה לא הספיקה. היה צריך \textit{לבנות} משהו שרץ על שתיהן — משהו שמכריח להבין את ההבדלים, לא רק לקרוא עליהם.

זה בדיוק מה שהפרויקט הזה עושה: מספק (provision) מכונת לינוקס חדשה והופך אותה לסביבה מאובטחת, מוגדרת ומוכנה לייצור — אוטומטית — גם על Ubuntu וגם על Fedora.

\subsection*{מה תלמד מהספר הזה?}

\begin{itemize}[noitemsep]
  \item מה זה לינוקס, מה זו הפצה, ומה ההבדל בין Ubuntu ל-Fedora
  \item כיצד כותבים סקריפטי Bash מקצועיים עם טיפול בשגיאות
  \item כיצד עובדים עם systemd — מנהל השירותים של לינוקס המודרני
  \item כיצד מנהלים משתמשים, הרשאות וחומות אש
  \item כיצד מתקינים ומגדירים nginx ו-Docker
  \item כיצד עובדים עם Git, GitHub Actions ו-Pull Requests
\end{itemize}

\begin{infobox}
\textbf{הערה:} הספר מיועד למי שמעולם לא עבד עם לינוקס, אבל גם מי שמכיר Ubuntu ימצא ערך רב בחלקים העוסקים בהשוואה בין ההפצות.
\end{infobox}

\newpage
%% ============================================================
\section{יסודות — מה זה לינוקס ומה ההבדל בין הפצות?}
%% ============================================================

\subsection{הגרעין (Kernel) וההפצה (Distribution)}

לינוקס עצמו הוא \textbf{גרעין} (Kernel) — תוכנה בסיסית שמנהלת את החומרה: מעבד, זיכרון, רשת, דיסק. אבל גרעין לבד הוא לא מספיק. צריך כלים, חבילות תוכנה, ממשק משתמש. שילוב של כל אלה יחד נקרא \textbf{הפצה} (Distribution).

\begin{center}
\begin{tikzpicture}[node distance=0.8cm]
  \node[block, fill=gray!20, minimum width=8cm] (kernel) {Linux Kernel};
  \node[block, fill=blue!15, minimum width=3.5cm, above left=of kernel] (ubuntu) {Ubuntu};
  \node[block, fill=red!15, minimum width=3.5cm, above right=of kernel] (fedora) {Fedora};
  \node[block, fill=purple!15, minimum width=3.5cm, above=1.6cm of kernel] (debian) {Debian / RHEL / Arch...};
  \draw[arrow] (ubuntu) -- (kernel);
  \draw[arrow] (fedora) -- (kernel);
  \draw[arrow] (debian) -- (kernel);
\end{tikzpicture}
\end{center}

\subsection{ההבדלים בין Ubuntu ו-Fedora}

\begin{center}
\begin{tabular}{lll}
\toprule
\textbf{נושא} & \textbf{Ubuntu} & \textbf{Fedora} \\
\midrule
בסיס & Debian & Red Hat \\
מנהל חבילות & \texttt{apt} & \texttt{dnf} \\
התקנת חבילה & \texttt{apt install nginx} & \texttt{dnf install nginx} \\
עדכון מערכת & \texttt{apt update \&\& apt upgrade} & \texttt{dnf update} \\
חומת אש & \texttt{ufw} & \texttt{firewalld} \\
פתיחת פורט & \texttt{ufw allow 80/tcp} & \texttt{firewall-cmd -{}-add-port=80/tcp} \\
מנהל שירותים & \texttt{systemctl} & \texttt{systemctl} (זהה!) \\
צביעת לוגים & \texttt{journalctl} & \texttt{journalctl} (זהה!) \\
\bottomrule
\end{tabular}
\end{center}

\begin{infobox}
\textbf{חשוב לזכור:} systemd ו-journalctl הם זהים בשתי ההפצות. ההבדלים הגדולים הם במנהל החבילות ובכלי חומת האש.
\end{infobox}

\subsection{מה זה Terminal ומה זה Shell?}

ה-\textbf{Terminal} הוא החלון שבו מקלידים פקודות. ה-\textbf{Shell} הוא התוכנה שמפרשת את הפקודות. ב-Linux, ה-Shell הנפוץ ביותר הוא \textbf{Bash}.

\begin{lstlisting}[language=bash, caption={דוגמה לפקודות בסיסיות ב-Bash}]
# ls — מציג את תוכן התיקייה הנוכחית
ls -la

# cd — מעבר לתיקייה
cd /etc/nginx

# cat — מציג את תוכן קובץ
cat /etc/os-release

# sudo — מריץ פקודה כ-root (מנהל מערכת)
sudo apt install nginx
\end{lstlisting}

\newpage
%% ============================================================
\section{כלים בסיסיים — Bash ו-systemd}
%% ============================================================

\subsection{קובץ .sh — מה זה?}

קובץ \texttt{.sh} הוא \textbf{סקריפט Bash} — קובץ טקסט פשוט שמכיל פקודות לינוקס, מבוצעות שורה אחר שורה. ההבדל היחיד בין הרצת פקודות בטרמינל לבין סקריפט הוא שבסקריפט הכל קורה אוטומטית.

\begin{lstlisting}[language=bash, caption={מבנה סקריפט בסיסי}]
#!/usr/bin/env bash
# השורה הראשונה נקראת "shebang" — מספרת ל-Linux איזה תוכנה תפרש את הקובץ

set -euo pipefail
# שלושה דגלי בטיחות:
# -e  : עצור מיד אם פקודה נכשלת
# -u  : שגיאה אם משתמשים במשתנה שלא הוגדר
# -o pipefail : אם פקודה ב-pipe נכשלת, כל ה-pipe נכשל

echo "Hello from bash!"
\end{lstlisting}

\subsection{משתנים ופונקציות}

\begin{lstlisting}[language=bash, caption={משתנים, פונקציות ו-if ב-Bash}]
# הגדרת משתנה — ללא רווחים סביב =
NAME="devops"

# שימוש במשתנה — עם $
echo "Hello, ${NAME}"

# פונקציה
log_info() {
    echo "[INFO] $*"
    # $* = כל הארגומנטים שהועברו לפונקציה
}

# if statement
if [[ "${EUID}" -ne 0 ]]; then
    log_info "חייבים להריץ כ-root"
    exit 1
fi
# EUID = מזהה המשתמש הנוכחי. root = 0
# -ne = "not equal" (השוואה מספרית)
\end{lstlisting}

\subsection{systemd — מנהל השירותים}

\textbf{systemd} הוא התוכנה שמנהלת את כל השירותים (processes) שרצים ברקע ב-Linux המודרני. כשהמחשב מתחיל, systemd הוא הדבר הראשון שרץ.

\begin{center}
\begin{tikzpicture}[node distance=1.2cm]
  \node[startend] (boot) {מחשב מופעל};
  \node[block, below=of boot] (kernel2) {Linux Kernel};
  \node[block, below=of kernel2] (systemd) {systemd (PID 1)};
  \node[block, below left=of systemd, xshift=-1cm] (nginx2) {nginx.service};
  \node[block, below=of systemd] (docker2) {docker.service};
  \node[block, below right=of systemd, xshift=1cm] (monitor) {nginx-monitor.service};

  \draw[arrow] (boot) -- (kernel2);
  \draw[arrow] (kernel2) -- (systemd);
  \draw[arrow] (systemd) -- (nginx2);
  \draw[arrow] (systemd) -- (docker2);
  \draw[arrow] (systemd) -- (monitor);
\end{tikzpicture}
\end{center}

\begin{center}
\begin{tabular}{ll}
\toprule
\textbf{פקודה} & \textbf{מה היא עושה} \\
\midrule
\texttt{systemctl start nginx} & מפעיל nginx עכשיו \\
\texttt{systemctl stop nginx} & עוצר nginx \\
\texttt{systemctl enable nginx} & מפעיל nginx בכל אתחול \\
\texttt{systemctl status nginx} & מציג את מצב nginx \\
\texttt{journalctl -u nginx} & מציג לוגים של nginx \\
\texttt{journalctl -u nginx -f} & מציג לוגים בזמן אמת \\
\bottomrule
\end{tabular}
\end{center}

\newpage
%% ============================================================
\section{מבנה הפרויקט}
%% ============================================================

\subsection{סקירה כללית}

\begin{lstlisting}[caption={מבנה תיקיות הפרויקט}]
multi-distro-provisioner/
|
|-- Vagrantfile          <- מגדיר את המכונות הוירטואליות
|-- Makefile             <- קיצורי דרך לפקודות
|-- CHECKLIST.md         <- מעקב אחר התקדמות
|
|-- docs/                <- תיעוד
|-- config/
|   |-- users.conf       <- רשימת המשתמשים ליצירה
|   +-- nginx/
|       +-- default.conf <- הגדרות שרת nginx
|
|-- scripts/
|   |-- common/          <- סקריפטים שרצים על שתי ההפצות
|   |-- ubuntu/          <- ספציפי ל-Ubuntu
|   +-- fedora/          <- ספציפי ל-Fedora
|
+-- systemd/             <- קבצי unit של systemd
\end{lstlisting}

\subsection{זרימת הפרויקט הכללית}

\begin{center}
\begin{tikzpicture}[node distance=0.9cm, every node/.style={font=\small}]
  \node[startend] (start) {\texttt{make up}};
  \node[block, below=of start] (vagrant) {Vagrant מפעיל VMs};
  \node[block, below=of vagrant] (provision) {Vagrant מריץ provision.sh};
  \node[block, below=of provision] (users) {שלב 1: יצירת משתמשים};
  \node[block, below=of users] (nginx3) {שלב 2: התקנת nginx};
  \node[block, below=of nginx3] (fw) {שלב 3: הגדרת חומת אש};
  \node[block, below=of fw] (docker3) {שלב 4: התקנת Docker};
  \node[block, below=of docker3] (sysd) {שלב 7: שירותי systemd};
  \node[startend, below=of sysd] (done) {מוכן!};

  \draw[arrow] (start) -- (vagrant);
  \draw[arrow] (vagrant) -- (provision);
  \draw[arrow] (provision) -- (users);
  \draw[arrow] (users) -- (nginx3);
  \draw[arrow] (nginx3) -- (fw);
  \draw[arrow] (fw) -- (docker3);
  \draw[arrow] (docker3) -- (sysd);
  \draw[arrow] (sysd) -- (done);
\end{tikzpicture}
\end{center}

\subsection{Vagrant — מה זה?}

\textbf{Vagrant} הוא כלי שמאפשר לתאר מכונה וירטואלית בקובץ טקסט (\texttt{Vagrantfile}) ולהפעיל אותה עם פקודה אחת. הרעיון הוא \textbf{Infrastructure as Code} — תשתית המתוארת בקוד.

\begin{lstlisting}[caption={קטע מה-Vagrantfile}]
Vagrant.configure("2") do |config|

  config.vm.define "ubuntu-node" do |ubuntu|
    ubuntu.vm.box     = "ubuntu/jammy64"  # Ubuntu 22.04
    ubuntu.vm.network "private_network", ip: "192.168.56.10"
    ubuntu.vm.network "forwarded_port", guest: 80, host: 8080
    ubuntu.vm.provision "shell",
      inline: "bash /vagrant/scripts/ubuntu/provision.sh"
  end

end
\end{lstlisting}

\newpage
%% ============================================================
\section{שלב 1 — ניהול משתמשים}
%% ============================================================

\subsection{מודל המשתמשים}

\begin{center}
\begin{tikzpicture}
  \node[block, fill=ubuntuorange!20, minimum width=4cm] (devops2) {devops};
  \node[block, fill=fedorablue!20, minimum width=4cm, right=2cm of devops2] (appuser2) {appuser};

  \node[below=0.4cm of devops2, font=\small, text=gray] (d1) {sudo: כן};
  \node[below=0.8cm of devops2, font=\small, text=gray] (d2) {docker group: כן};
  \node[below=1.2cm of devops2, font=\small, text=gray] (d3) {/home/devops};

  \node[below=0.4cm of appuser2, font=\small, text=gray] (a1) {sudo: לא};
  \node[below=0.8cm of appuser2, font=\small, text=gray] (a2) {docker group: כן};
  \node[below=1.2cm of appuser2, font=\small, text=gray] (a3) {בעלים של /app};
\end{tikzpicture}
\end{center}

\subsection{זרימת יצירת משתמש}

\begin{center}
\begin{tikzpicture}[node distance=0.8cm]
  \node[startend] (s) {קרא שורה מ-users.conf};
  \node[decision, below=of s] (grp) {קבוצה קיימת?};
  \node[block, right=2cm of grp] (mkgrp) {\texttt{groupadd}};
  \node[decision, below=1.2cm of grp] (usr) {משתמש קיים?};
  \node[block, right=2cm of usr] (mkusr) {\texttt{useradd}};
  \node[decision, below=1.2cm of usr] (sud) {sudo=yes?};
  \node[block, right=2cm of sud] (mksud) {כתוב ל-sudoers.d};
  \node[startend, below=1.2cm of sud] (next) {משתמש הבא};

  \draw[arrow] (s) -- (grp);
  \draw[arrow] (grp) -- node[above,font=\tiny]{לא} (mkgrp);
  \draw[arrow] (grp) -- node[left,font=\tiny]{כן} (usr);
  \draw[arrow] (mkgrp) |- (usr);
  \draw[arrow] (usr) -- node[above,font=\tiny]{לא} (mkusr);
  \draw[arrow] (usr) -- node[left,font=\tiny]{כן} (sud);
  \draw[arrow] (mkusr) |- (sud);
  \draw[arrow] (sud) -- node[above,font=\tiny]{כן} (mksud);
  \draw[arrow] (sud) -- node[left,font=\tiny]{לא} (next);
  \draw[arrow] (mksud) |- (next);
\end{tikzpicture}
\end{center}

\subsection{הסבר הקוד}

\begin{lstlisting}[language=bash, caption={create\_users.sh — חלק מרכזי}]
# mapfile קורא את כל הקובץ למערך לפני הלולאה
# זה מונע בעיות stdin כשפקודות כמו useradd רצות בתוך הלולאה
mapfile -t lines < "${CONFIG_FILE}"

for line in "${lines[@]}"; do
    # דלג על שורות ריקות ושורות הערה
    [[ -z "${line}" || "${line}" =~ ^# ]] && continue

    # פרסר את השורה לפי ":"
    # <<< מזין את המחרוזת ישירות ל-read
    IFS=: read -r username group shell sudo_access <<< "${line}"

    # צור קבוצה אם לא קיימת
    if ! getent group "${group}" &>/dev/null; then
        groupadd "${group}"
    fi

    # צור משתמש אם לא קיים
    if ! id "${username}" &>/dev/null; then
        useradd --create-home --gid "${group}" \
                --shell "${shell}" "${username}"
    fi

    # הגדר sudo אם נדרש
    if [[ "${sudo_access}" == "yes" ]]; then
        echo "${username} ALL=(ALL) NOPASSWD:ALL" \
             > "/etc/sudoers.d/${username}"
        chmod 440 "/etc/sudoers.d/${username}"
        # chmod 440 חובה — sudo מסרב לקרוא קבצים עם הרשאות שגויות
    fi
done
\end{lstlisting}

\newpage
%% ============================================================
\section{שלב 2 — nginx}
%% ============================================================

\subsection{מה זה nginx?}

\textbf{nginx} הוא שרת web — תוכנה שמאזינה לבקשות HTTP ומחזירה תגובות. כשנכתוב \texttt{http://localhost:8080} בדפדפן, nginx הוא זה שמחזיר את דף ה-HTML.

\subsection{כיצד nginx עובד עם systemd}

\begin{center}
\begin{tikzpicture}[node distance=1cm]
  \node[block] (enable) {\texttt{systemctl enable nginx}};
  \node[block, right=1.5cm of enable] (start2) {\texttt{systemctl start nginx}};
  \node[block, below=of enable] (symlink) {יוצר symlink ב-\texttt{/etc/systemd/system/}};
  \node[block, below=of start2] (running) {nginx רץ עכשיו};
  \node[block, below=1.5cm of symlink] (reboot) {אתחול הבא: nginx עולה אוטומטי};

  \draw[arrow] (enable) -- (symlink);
  \draw[arrow] (start2) -- (running);
  \draw[arrow] (symlink) -- (reboot);
\end{tikzpicture}
\end{center}

\begin{lstlisting}[language=bash, caption={install\_nginx.sh — התקנה והפעלה}]
# זיהוי ההפצה — אוטומטי דרך /etc/os-release
distro="$(detect_distro)"

case "${distro}" in
    ubuntu|debian)
        apt-get update -qq        # רענון רשימת החבילות
        apt-get install -y nginx  # -y = כן לכל השאלות
        ;;
    fedora|rhel)
        dnf install -y nginx      # dnf לא צריך update נפרד
        ;;
esac

# העתק את קובץ הגדרות nginx שלנו
cp "${CONFIG_FILE}" /etc/nginx/conf.d/default.conf

# nginx -t מריץ dry-run של parser הקונפיגורציה
# אם יש שגיאת syntax, הסקריפט נעצר כאן (בגלל set -e)
nginx -t

systemctl enable nginx   # הפעל בכל אתחול
systemctl start nginx    # הפעל עכשיו
\end{lstlisting}

\subsection{קובץ הגדרות nginx}

\begin{lstlisting}[caption={config/nginx/default.conf}]
server {
    listen 80;          # האזן על פורט 80 (HTTP)
    server_name _;      # _ = קבל כל hostname

    root /var/www/html; # תיקיית קבצי האתר
    index index.html;

    location / {
        try_files $uri $uri/ =404;
    }

    # endpoint לבדיקת תקינות — משמש את סקריפט הניטור
    location /health {
        return 200 "OK\n";
        add_header Content-Type text/plain;
    }
}
\end{lstlisting}

\newpage
%% ============================================================
\section{שלב 3 — חומת אש (Firewall)}
%% ============================================================

\subsection{מה זו חומת אש?}

חומת אש היא תוכנה שמחליטה אילו חיבורי רשת מותרים. ברירת המחדל הבטוחה היא: \textbf{חסום הכל, אלא אם כן אישרנו במפורש}.

\subsection{ההבדל בין ufw ל-firewalld}

\begin{center}
\begin{tabular}{lll}
\toprule
\textbf{פעולה} & \textbf{Ubuntu (ufw)} & \textbf{Fedora (firewalld)} \\
\midrule
הפעל & \texttt{ufw enable} & \texttt{systemctl start firewalld} \\
חסום הכל & \texttt{ufw default deny incoming} & ברירת מחדל \\
אפשר SSH & \texttt{ufw allow ssh} & \texttt{firewall-cmd -{}-add-service=ssh} \\
אפשר HTTP & \texttt{ufw allow 80/tcp} & \texttt{firewall-cmd -{}-add-service=http} \\
שמור לצמיתות & אוטומטי & \texttt{-{}-permanent + -{}-reload} \\
הצג חוקים & \texttt{ufw status numbered} & \texttt{firewall-cmd -{}-list-all} \\
\bottomrule
\end{tabular}
\end{center}

\begin{infobox}
\textbf{המלכודת הכי נפוצה ב-Fedora:} \texttt{-{}-permanent} שומר את החוק בדיסק אבל לא מפעיל אותו. חייבים לקרוא \texttt{firewall-cmd -{}-reload} אחרי כל שינוי!
\end{infobox}

\begin{lstlisting}[language=bash, caption={setup\_firewall.sh — Ubuntu}]
# חשוב: הגדר ברירות מחדל לפני הפעלת ufw
# אחרת עלול לנעול את עצמך מחוץ ל-SSH!
ufw default deny incoming   # חסום הכל
ufw default allow outgoing  # אפשר יציאה

# SSH חייב להיות ראשון — אחרת נאבד גישה
ufw allow ssh               # = ufw allow 22/tcp

ufw allow 80/tcp            # HTTP
ufw allow 443/tcp           # HTTPS

ufw --force enable          # --force מונע שאלת אישור אינטראקטיבית
ufw reload
\end{lstlisting}

\newpage
%% ============================================================
\section{שלב 4 — Docker}
%% ============================================================

\subsection{ארכיטקטורת Docker}

\begin{center}
\begin{tikzpicture}[node distance=1.2cm]
  \node[block] (cli) {\texttt{docker} CLI};
  \node[block, below=of cli] (socket) {\texttt{/var/run/docker.sock} (Unix Socket)};
  \node[block, below=of socket] (daemon) {\texttt{dockerd} (daemon)};
  \node[block, below=of daemon] (containerd) {\texttt{containerd} (container runtime)};

  \draw[arrow] (cli) -- node[right,font=\tiny]{פקודות} (socket);
  \draw[arrow] (socket) -- (daemon);
  \draw[arrow] (daemon) -- (containerd);

  \node[right=2cm of daemon, block, fill=green!15] (c1) {Container 1: nginx};
  \node[right=2cm of containerd, block, fill=green!15] (c2) {Container 2: postgres};
  \draw[arrow] (containerd) -- (c1);
  \draw[arrow] (containerd) -- (c2);
\end{tikzpicture}
\end{center}

\subsection{למה צריך GPG Key?}

כשמוסיפים repository של Docker, Linux צריך לדעת שהחבילות הן אמיתיות ולא זויפו. Docker חותמים על החבילות שלהם עם \textbf{מפתח פרטי} (GPG private key). אנחנו מורידים את ה\textbf{מפתח ציבורי} ושומרים אותו — ו-apt/dnf יוודאו כל חבילה מולו.

\begin{lstlisting}[language=bash, caption={setup\_docker.sh — הוספת GPG key ו-repository}]
# הורד את ה-GPG key של Docker ושמור אותו
curl -fsSL https://download.docker.com/linux/ubuntu/gpg \
    | gpg --dearmor -o /etc/apt/keyrings/docker.gpg

# הוסף את ה-repository של Docker ל-apt
# arch=$(dpkg --print-architecture) = amd64
# VERSION_CODENAME = jammy (Ubuntu 22.04)
echo "deb [arch=$(dpkg --print-architecture) \
    signed-by=/etc/apt/keyrings/docker.gpg] \
    https://download.docker.com/linux/ubuntu \
    $(. /etc/os-release && echo "${VERSION_CODENAME}") stable" \
    | tee /etc/apt/sources.list.d/docker.list

apt-get update -qq
apt-get install -y docker-ce docker-ce-cli containerd.io
\end{lstlisting}

\newpage
%% ============================================================
\section{שלב 5 — סקריפטי Bash}
%% ============================================================

\subsection{health\_check.sh — בדיקת תקינות}

\begin{lstlisting}[language=bash, caption={health\_check.sh — מערך שגיאות ופונקציית בדיקה}]
# מערך לאיסוף כשלונות — נבדוק הכל ונדווח בסוף
FAILURES=()

run_check() {
    local name="$1"
    local cmd="$2"

    if eval "${cmd}" &>/dev/null; then
        echo "  OK  ${name}"
    else
        echo "  FAIL  ${name}"
        FAILURES+=("${name}")   # += מוסיף לסוף המערך
    fi
}

run_check "nginx רץ"    "systemctl is-active nginx"
run_check "nginx תקין"  "nginx -t"
run_check "HTTP עובד"   "curl -sf http://localhost/health"

# ${#FAILURES[@]} = אורך המערך
if [[ "${#FAILURES[@]}" -eq 0 ]]; then
    echo "הכל תקין!"
    exit 0
else
    echo "${#FAILURES[@]} בדיקות נכשלו"
    exit 1
fi
\end{lstlisting}

\subsection{user\_report.sh — דוח משתמשים}

\begin{lstlisting}[language=bash, caption={user\_report.sh — שימוש ב-awk ו-printf}]
# printf לעיצוב עמודות מיושרות
# %-15s = מחרוזת, מיושרת שמאלה, רוחב 15 תווים
printf "%-15s %-20s %-10s\n" "משתמש" "קבוצות" "sudo"

# /etc/passwd פורמט: username:x:UID:GID:comment:home:shell
# UID >= 1000 = משתמשים אנושיים (לא מערכת)
while IFS=: read -r username shell; do
    user_groups=$(groups "${username}" | cut -d: -f2)

    if [[ -f "/etc/sudoers.d/${username}" ]]; then
        sudo_status="כן"
    else
        sudo_status="לא"
    fi

    printf "%-15s %-20s %-10s\n" \
        "${username}" "${user_groups}" "${sudo_status}"

done < <(awk -F: '$3 >= 1000 {print $1 ":" $7}' /etc/passwd)
\end{lstlisting}

\newpage
%% ============================================================
\section{שלב 6 — ניטור לוגים}
%% ============================================================

\subsection{מה זה journalctl?}

כל שירות שמנוהל על ידי systemd כותב את הלוגים שלו ל-\textbf{journal} — מאגר לוגים מרכזי של המערכת. \texttt{journalctl} הוא הכלי לקרוא אותם.

\begin{center}
\begin{tabular}{ll}
\toprule
\textbf{פקודה} & \textbf{מה היא עושה} \\
\midrule
\texttt{journalctl -u nginx} & כל הלוגים של nginx \\
\texttt{journalctl -u nginx -f} & לוגים בזמן אמת \\
\texttt{journalctl -u nginx --since "1 hour ago"} & שעה אחורה \\
\texttt{journalctl -u nginx -n 20} & 20 שורות אחרונות \\
\texttt{journalctl -u nginx -p err} & רק שגיאות \\
\bottomrule
\end{tabular}
\end{center}

\subsection{קודי סטטוס HTTP}

\begin{center}
\begin{tabular}{lll}
\toprule
\textbf{טווח} & \textbf{משמעות} & \textbf{דוגמה} \\
\midrule
2xx & הצלחה & 200 OK \\
3xx & הפניה & 301 Moved Permanently \\
4xx & שגיאת לקוח & 404 Not Found \\
5xx & שגיאת שרת & 500 Internal Server Error \\
\bottomrule
\end{tabular}
\end{center}

\begin{lstlisting}[language=bash, caption={monitor\_logs.sh — ספירת שגיאות 5xx}]
LOOKBACK="1 hour ago"
ERROR_THRESHOLD=10

# קרא לוגים מהשעה האחרונה
NGINX_LOGS=$(journalctl -u nginx \
    --since "${LOOKBACK}" \
    --no-pager -o short 2>/dev/null || true)

# grep -cE = ספור שורות תואמות לביטוי רגולרי
# ' 5[0-9]{2} ' = פורט HTTP שמתחיל ב-5 (5xx)
ERROR_COUNT=$(echo "${NGINX_LOGS}" \
    | grep -cE ' 5[0-9]{2} ' || true)

if [[ "${ERROR_COUNT}" -gt "${ERROR_THRESHOLD}" ]]; then
    echo "ALERT: ${ERROR_COUNT} שגיאות 5xx בשעה האחרונה!"
    exit 2
fi
\end{lstlisting}

\newpage
%% ============================================================
\section{שלב 7 — שירותי systemd}
%% ============================================================

\subsection{מבנה קובץ .service}

\begin{lstlisting}[caption={nginx-monitor.service — הסבר מלא}]
[Unit]
# תיאור השירות
Description=nginx health monitor

# התחל רק אחרי ש-nginx עלה
After=nginx.service network.target

[Service]
# Type=simple: systemd מניח שהשירות רץ ישירות (לא fork לרקע)
Type=simple

# הפקודה להרצה — חייב נתיב מלא
ExecStart=/usr/local/bin/nginx-monitor-loop.sh

# Restart=always: הפעל מחדש אם עצר מכל סיבה
Restart=always
RestartSec=10

# לוגים יגיעו ל-journal (ייראו ב-journalctl -u nginx-monitor)
StandardOutput=journal
StandardError=journal

[Install]
# WantedBy=multi-user.target: הפעל בסביבה רגילה (לאחר boot)
WantedBy=multi-user.target
\end{lstlisting}

\subsection{Timer — חלופה מודרנית ל-cron}

\begin{center}
\begin{tabular}{ll}
\toprule
\textbf{cron (ישן)} & \textbf{systemd timer (מודרני)} \\
\midrule
\texttt{0 * * * * /script.sh} & \texttt{OnUnitActiveSec=1h} \\
לוגים? לא \\
\quad & לוגים ב-journal \\
דרוש cron daemon & מנוהל על ידי systemd \\
הרצה ידנית? מסובך & \texttt{systemctl start log-alert} \\
\bottomrule
\end{tabular}
\end{center}

\begin{lstlisting}[caption={log-alert.timer}]
[Timer]
# הרץ 5 דקות אחרי אתחול
OnBootSec=5min

# ואז כל שעה
OnUnitActiveSec=1h

# אם המחשב היה כבוי בזמן שהיה צריך לרוץ — הרץ מיד באתחול
Persistent=true

[Install]
WantedBy=timers.target
\end{lstlisting}

\subsection{פקודות ניהול השירותים}

\begin{lstlisting}[language=bash]
# לאחר העתקת קבצי .service — חובה!
systemctl daemon-reload

# הפעל והגדר להפעלה אוטומטית
systemctl enable --now nginx-monitor.service

# הפעל את ה-timer
systemctl enable --now log-alert.timer

# ראה מצב השירות
systemctl status nginx-monitor

# ראה לוגים חיים
journalctl -u nginx-monitor -f

# ראה כל ה-timers
systemctl list-timers
\end{lstlisting}

\newpage
%% ============================================================
\section{Git ו-GitHub — עבודה מקצועית}
%% ============================================================

\subsection{זרימת עבודה עם Git}

\begin{center}
\begin{tikzpicture}[node distance=1cm]
  \node[block, fill=green!20, minimum width=5cm] (main2) {\texttt{main} (branch יציב)};
  \node[block, below left=1.5cm and 0.5cm of main2, minimum width=4cm] (feat) {\texttt{feat/phase-1-users}};
  \node[block, below right=1.5cm and 0.5cm of main2, minimum width=4cm] (feat2) {\texttt{feat/phase-2-nginx}};

  \draw[arrow] (main2) -- node[left,font=\tiny]{\texttt{checkout -b}} (feat);
  \draw[arrow] (main2) -- node[right,font=\tiny]{\texttt{checkout -b}} (feat2);
  \draw[arrow, bend left=30] (feat) to node[left,font=\tiny]{Pull Request} (main2);
  \draw[arrow, bend right=30] (feat2) to node[right,font=\tiny]{Pull Request} (main2);
\end{tikzpicture}
\end{center}

\begin{lstlisting}[language=bash, caption={זרימת עבודה לכל שלב}]
# 1. עבור ל-main ועדכן
git checkout main
git pull

# 2. צור branch חדש לשלב הנוכחי
git checkout -b feat/phase-1-users

# 3. עבוד, ערוך קבצים, בדוק

# 4. הוסף לקומיט
git add scripts/common/create_users.sh
git commit -m "feat(users): implement user creation from config file"

# 5. שלח ל-GitHub
git push --set-upstream origin feat/phase-1-users

# 6. פתח Pull Request ב-GitHub עם: Closes #2
\end{lstlisting}

\subsection{GitHub Actions — CI אוטומטי}

\begin{center}
\begin{tikzpicture}[node distance=1cm]
  \node[startend] (push2) {git push};
  \node[block, below=of push2] (trigger) {GitHub Actions מופעל};
  \node[block, below=of trigger] (checkout2) {\texttt{actions/checkout@v4}};
  \node[block, below=of checkout2] (shell) {shellcheck על כל קבצי .sh};
  \node[decision, below=of shell] (ok) {עבר?};
  \node[block, right=2cm of ok, fill=green!20] (merge) {PR מאושר למיזוג};
  \node[block, left=2cm of ok, fill=red!20] (block) {PR חסום};

  \draw[arrow] (push2) -- (trigger);
  \draw[arrow] (trigger) -- (checkout2);
  \draw[arrow] (checkout2) -- (shell);
  \draw[arrow] (shell) -- (ok);
  \draw[arrow] (ok) -- node[above,font=\tiny]{כן} (merge);
  \draw[arrow] (ok) -- node[above,font=\tiny]{לא} (block);
\end{tikzpicture}
\end{center}

\newpage
%% ============================================================
\section{הרצת הפרויקט — מדריך מלא}
%% ============================================================

\subsection{דרישות מוקדמות}

\begin{enumerate}[noitemsep]
  \item \textbf{VirtualBox} — מנוע המכונות הוירטואליות: \url{https://www.virtualbox.org}
  \item \textbf{Vagrant} — מנהל המכונות: \url{https://developer.hashicorp.com/vagrant}
  \item \textbf{make} (Windows): \texttt{winget install GnuWin32.Make}
\end{enumerate}

\subsection{הפעלה ראשונה}

\begin{lstlisting}[language=bash, caption={הרצת הפרויקט לראשונה}]
# שלב 1: הפעל את שתי המכונות
# בפעם הראשונה יוריד תמונות (~1GB) — ייקח 5-10 דקות
make up

# שלב 2: בדוק שnginx עובד מהדפדפן
# Ubuntu: http://localhost:8080
# Fedora: http://localhost:8081

# שלב 3: כנס ל-Ubuntu ובדוק
make ssh-ubuntu

# בתוך המכונה:
sudo bash /opt/provisioner/scripts/common/health_check.sh
sudo bash /opt/provisioner/scripts/common/user_report.sh
sudo ufw status numbered
systemctl status nginx docker nginx-monitor
journalctl -u nginx-monitor --no-pager -n 20
exit

# שלב 4: בדוק idempotency — הרץ שוב, אמור לעבוד ללא שגיאות
make provision-ubuntu
\end{lstlisting}

\subsection{פקודות שימושיות}

\begin{center}
\begin{tabular}{ll}
\toprule
\textbf{פקודה} & \textbf{מה היא עושה} \\
\midrule
\texttt{make up} & הפעל שתי המכונות \\
\texttt{make provision-ubuntu} & הרץ provisioner על Ubuntu \\
\texttt{make provision-fedora} & הרץ provisioner על Fedora \\
\texttt{make ssh-ubuntu} & פתח shell ב-Ubuntu \\
\texttt{make ssh-fedora} & פתח shell ב-Fedora \\
\texttt{make down} & כבה המכונות \\
\texttt{make clean} & מחק המכונות לגמרי \\
\bottomrule
\end{tabular}
\end{center}

%% ============================================================
\section{סיכום}
%% ============================================================

פרויקט זה לימד אותי שהדרך הטובה ביותר ללמוד לינוקס היא לבנות משהו שצריך לרוץ עליו — על הפצות שונות, בסביבות שונות.

ההבדלים בין Ubuntu ו-Fedora אינם רק בשמות הפקודות. הם מייצגים פילוסופיות שונות של ניהול מערכת: Ubuntu מכוונת לנגישות ומהירות התקנה, Fedora מכוונת לחדשנות ולבסיס לסביבות ארגוניות (RHEL). הבנת שניהם נותנת תמונה מלאה של עולם לינוקס.

הכלים שנלמדו כאן — systemd, bash scripting, firewall management, Docker, CI/CD עם GitHub Actions — הם הכלים שמשמשים מהנדסי DevOps ומנהלי מערכת ב-RedHat, AWS, Google, ו-Meta בכל יום.

\vspace{1cm}
\begin{infobox}
\centering
\textbf{הקוד המלא זמין ב-GitHub:}\\
\url{https://github.com/EsteeCohen/multi-distro-provisioner}
\end{infobox}

\end{document}